% Pista ana-25s-q1 · Anàlisi
% Procedència: PAU Catalunya 2025 · Convocatòria de setembre · Sèrie 3 · Exercici 1
\documentclass[11pt,a4paper]{article}
\usepackage{pista}
\pistainfo{Catalunya 2025}{Convocatòria de setembre}{Sèrie 3}

\begin{document}

\section*{\textcolor{blauPista}{Exercici 1}}

\noindent Una família vol comprar un terreny per a fer-s'hi una casa envoltada de penya-segats amb vistes al mar. En aquella zona de la costa, els penya-segats segueixen les rectes $y = 0$ i $y = 3x$. A més, la família vol que el terreny sigui triangular i que el tercer costat del triangle passi pel punt $P = (1, 1)$, tal com es pot veure en la figura.

\noindent\textit{(A la figura: el terreny és el triangle del primer quadrant limitat per l'eix $OX$, la recta $y = 3x$ i la recta $r$, que passa per $P$ i és decreixent.)}

\subsection*{\textit{a)} \normalfont Plantegeu l'equació de la recta $r$ que defineix el tercer costat del triangle en funció del seu pendent $m$, i comproveu que l'àrea del terreny ve donada per
\[
A(m) = \dfrac{3}{2} \cdot \dfrac{m^{2} - 2m + 1}{m^{2} - 3m}.
\]
\hspace*{\fill}\textnormal{[1~punt]}}

\pista{En primer lloc, escriu la recta $r$ en la forma punt-pendent:
\[
y - y_{0} = m\,(x - x_{0}),
\]
on $(x_{0}, y_{0})$ és un punt de la recta, en aquest cas $P$, i $m$ és el pendent.

\medskip
A continuació, fixa't a la figura que la base del triangle és sobre l'eix d'abscisses, des de l'origen fins al punt on $r$ talla aquest eix: per trobar-lo, fes $y = 0$ a l'equació d'$r$. L'altura és la segona coordenada del punt on $r$ talla la recta $y = 3x$: per trobar-lo, substitueix $y = 3x$ a l'equació d'$r$. Tots dos punts et quedaran en funció de $m$. No és tan difícil com sembla, perquè cada vegada només has d'aïllar la $x$ d'una equació de primer grau.

\medskip
Finalment, calcula l'àrea i simplifica-la fins a arribar a l'expressió de l'enunciat:
\[
\text{Àrea} = \frac{\text{base} \cdot \text{altura}}{2}
\]}

\subsection*{\textit{b)} \normalfont Calculeu el valor de $m$ que fa que l'àrea d'aquest terreny (i, per tant, el seu preu) sigui mínima. Quin és el valor d'aquesta àrea? \hspace*{\fill}\textnormal{[1,5~punts]}}

\pista{En primer lloc, deriva $A(m)$ amb la regla de la derivada d'un quocient:
\[
\left(\frac{u}{v}\right)' = \frac{u' \cdot v - u \cdot v'}{v^{2}}
\]
El $\frac{3}{2}$ multiplica tota la fracció: el pots deixar fora i multiplicar-lo al final. A continuació, resol $A'(m) = 0$: només cal igualar a zero el numerador. Obtindràs dues solucions. Fixa't a la figura que la recta $r$ és decreixent, i per tant el seu pendent ha de ser negatiu: queda't amb la solució que compleix aquesta condició.

\medskip
Després, comprova que efectivament has obtingut un mínim: pots fer-ho estudiant el signe de la funció $A'$. Com que tens una funció racional i el denominador en aquest cas sempre serà positiu, només cal que estudiïs el signe del numerador.

\medskip
Finalment, tingues en compte que la teva resposta ha d'incloure, d'una banda, el valor de $m$; i d'altra banda, l'àrea mínima.}

\end{document}
